\documentclass[12pt,a4paper]{article}
\usepackage[utf8]{inputenc}
\usepackage{amsmath}
\usepackage{amsfonts}
\usepackage{amssymb}
\usepackage{geometry}
\usepackage{booktabs}
\usepackage{hyperref}

\geometry{margin=1in}

\begin{document}

\begin{center}
    \large \textbf{The Principle of Planck’s Field: Resolving the Dark Matter Paradox via Density Satiation and the Metric Impedance Scalar} \\
    \vspace{1em}
    \normalsize \textbf{Robin Skavberg} \\
    \normalsize Independent Researcher \\
    \normalsize ORCID: 0009-0003-9126-6196 \\
    \normalsize Email: skavberg@gmail.com \\
    \vspace{1em}
    \normalsize February 2026
\end{center}

\begin{abstract}
This paper formalizes the "Frozen Core" identity of Dark Matter by demonstrating that the Schwarzschild boundary is a density-dependent phase transition rather than a mass-dependent spatial coordinate. By applying the variable gravitational resultant $G(\rho_m)$, we prove that matter reaching the field satiation threshold $\rho_{sat}$ loses all vibrational degrees of freedom, rendering it electromagnetically silent and thermodynamically at absolute zero. We validate this model by demonstrating that the neutrino naturally exists in this satiated state, identifying it as the fundamental subatomic constituent of Dark Matter.
\end{abstract}

\section{Introduction: Foundational Field Axioms}
As established in previous derivations, we treat the gravitational constant $G$ not as a universal constant, but as the restorative response of the Planck Field to matter-energy density. This relationship is defined by:
\begin{equation}
G(\rho_m) = \frac{c^2}{8\pi\rho_m}
\end{equation}
Where $c^2$ represents the field's restorative stiffness. Under this framework, the Schwarzschild boundary condition ($R_s = 2GM/c^2$) undergoes a critical transformation. Substituting Eq.(1) into the Schwarzschild formula:
\begin{equation}
R_s = \frac{2 \left( \frac{c^2}{8\pi\rho_m} \right) M}{c^2} = \frac{M}{4\pi\rho_m}
\end{equation}
Given that $M = \rho_m \cdot \frac{4}{3}\pi R^3$ for a spherical volume, substitution yields:
\begin{equation}
R = \frac{\rho_m \cdot \frac{4}{3}\pi R^3}{4\pi\rho_m} \implies R = \frac{1}{3}R^3
\end{equation}
This identity confirms that at the satiation limit, the boundary is a geometric constant of the field's displacement capacity, independent of the specific mass $M$.

\section{The Frozen Core as Dark Matter}
We define Dark Matter as matter that has reached the density satiation threshold $\rho_{sat} \approx 5.358 \times 10^{25}$ kg/m$^3$. At this limit, the field's capacity for displacement is exhausted.

\subsection{The Metric Impedance Scalar ($\zeta$)}
The field’s ability to support wave-like perturbations (vibrations) is governed by the Metric Impedance Scalar:
\begin{equation}
\zeta = \frac{1}{1 - (\rho_m / \rho_{sat})}
\end{equation}
As $\rho_m \to \rho_{sat}$, $\zeta \to \infty$. This infinite impedance creates a "Frozen Core" where internal degrees of freedom are suppressed to zero.

\subsection{Electromagnetic and Thermal Silence}
Thermal energy is the manifestation of field vibrations ($E_{th} \propto k_B T$). In a satiated state, the lack of vibrational "slack" in the medium results in:
\begin{equation}
\Delta \text{Vibration} \to 0 \implies T \to 0 \text{ K}
\end{equation}
Without vibrational degrees of freedom, the matter cannot emit, absorb, or reflect light, accounting for the "dark" profile observed in astrophysics while maintaining its gravitational signature.

\section{Scale Invariance of the Satiated State}
A critical requirement for Dark Matter is the ability to exist across various scales. We demonstrate that the Frozen Core identity is strictly density-dependent. 

Consider a mass $M$ at $\rho_{sat}$. If the mass is scaled to $M' = kM$, the volume must scale to $V' = kV$ to maintain the state. Because $\zeta$ is a function of density ratio, the "Dark" state persists whether the radius is subatomic or galactic. This proves that Dark Matter is not a specific "species" of particle, but a physical phase of matter.

\section{Evaluation of the Neutrino}
We apply this density-satiation model to the neutrino. Using 2026 benchmark measurements:
\begin{itemize}
    \item \textbf{Mass ($m_\nu$):} $\approx 8 \times 10^{-37}$ kg
    \item \textbf{Effective Radius ($r_\nu$):} $\approx 10^{-21}$ m (based on weak interaction limits)
\end{itemize}
The resulting density $\rho_\nu$ is:
\begin{equation}
\rho_\nu = \frac{m_\nu}{\frac{4}{3}\pi r_\nu^3} \approx 1.9 \times 10^{26} \text{ kg/m}^3
\end{equation}
Comparing this to our calculated field satiation limit ($\rho_{sat} \approx 5.358 \times 10^{25}$ kg/m$^3$), we find that $\rho_\nu > \rho_{sat}$. This classifies the neutrino as a fundamental subatomic Frozen Core, naturally explaining its "dark" properties and minimal interaction cross-section.

\section{Resolution of Field Pressure and Observations}
The interaction between Dark Matter and Baryonic matter is mediated by the gradient of the Field's Scalar Pressure ($P_\Phi$).

\subsection{The Cusp-Core Resolution}
Standard Cold Dark Matter (CDM) models predict a density "cusp" at galactic centers. However, the Principle of Planck’s Field dictates that once density reaches $\rho_{sat}$, the field becomes structurally rigid. This rigidity acts as a "hard floor," preventing infinite compression and naturally resulting in the flat-density "cores" observed in spiral galaxies.

\section{Conclusion}
By redefining $G$ as a field resultant, we reconcile the Dark Matter paradox. Dark Matter is simply matter—predominantly neutrinos—that has reached the field's maximum displacement limit. These Frozen Cores are non-vibrational, thermodynamically at absolute zero, and scale-invariant, providing a unified mechanical explanation for both subatomic particles and black holes.

\begin{thebibliography}{9}
\bibitem{Skavberg2026a} Skavberg, R. (2026). \textit{A Mechanical Route to G: Matching Dark-Matter Field Pressure to the Einstein-Hilbert Coupling}. Zenodo.
\bibitem{Skavberg2026b} Skavberg, R. (2026). \textit{A Mechanical Extension of General Relativity: The Singularity in Equilibrium}. Zenodo.
\end{thebibliography}

\end{document}